\label{mab:rm:ch:chaos}

\section{Construction intrinsèque de la famille dodécaphasique}

Les constantes de Feigenbaum caractérisent l’universalité du
doublement de période. Leur dérivation dans HMT exige de construire au préalable
une famille analytique unimodale, de démontrer sa criticité quadratique et
d’étudier ensuite l’action de l’opérateur de renormalisation.

L’application initiale ne procède pas d’un choix harmonique externe. Le retour
de \(\TPK\) détermine la roue dodécaphasique orientée
\(U_{12}^{\rm sign}\). Dans le secteur uniforme \(\eta=0\), la trace
normalisée de ses douze secteurs fixe \(w_m=1/12\), et la fonctionnelle de phase
\begin{equation}
\mathcal P_{\rm crit}:U_{12}^{\rm sign}\longrightarrow\R/2\pi\Z,
\qquad
m\longmapsto \frac{(30m-10)\pi}{180}
=\pi\frac{3m-1}{18}
\label{mab:rm:eq:critical-phase-reader}
\end{equation}
conserve l’ordre et l’orientation. La composition causale est
\[
\APP\to\TRIT\to\TPK\to U_{12}^{\rm sign}
\xrightarrow{\mathcal P_{\rm crit}}
\{(\phi_m,w_m)\}_{m=1}^{12}
\to u_0\to f_\lambda.
\]
Toute modification de \(\mathcal P_{\rm crit}\), des poids ou de l’ordre des
phases altère la famille et constitue un critère de réfutation. Ces
choix doivent être fixés avant la confrontation avec \(\delta_F\).

La famille dodécaphasique résultante part de

\[
a_m=\frac{3m-1}{18},\qquad
\phi_m=\pi a_m,\qquad m=1,\ldots,12,
\]

et du deuxième moment

\[
D_0=\frac1{12}\sum_{m=1}^{12}a_m^2.
\]

L’identité

\begin{equation}
\sum_{m=1}^{12}(3m-1)^2=5394=6\cdot899
\label{mab:rm:eq:chaos-sum2}
\end{equation}

produit

\[
D_0=\frac{899}{648},\qquad
s_0=\frac{36}{\pi\sqrt{899}}.
\]

Par conséquent,

\begin{equation}
s_0\phi_m=\frac{2(3m-1)}{\sqrt{899}}.
\label{mab:rm:eq:pi-cancellation}
\end{equation}

Le facteur \(\pi\) s’annule avant d’effectuer les itérations. La famille ne
dépend ni de \(\delta_F\) ni de \(\alpha_F\) comme données d’entrée.

L’annulation est une propriété structurelle de la normalisation. Avant de
normaliser,
le deuxième moment angulaire est
\[
\frac1{12}\sum_{m=1}^{12}\phi_m^2
=\pi^2D_0.
\]
La condition de criticité quadratique fixe
\(s_0^2\pi^2D_0=2\), et chaque phase normalisée dépend donc seulement de
l’arithmétique entière de \(3m-1\). La fonctionnelle conserve la provenance angulaire de
la roue, mais l’itération reçoit une famille réelle intrinsèque. C’est
précisément ce qui permet de distinguer une dérivation d’un étalonnage
par \(\pi\).

\section{Représentation analytique exacte et expression trigonométrique fermée}

On définit

\begin{equation}
u_0(x)=1-\frac1{12}\sum_{m=1}^{12}
\cos\!\left(\frac{2(3m-1)}{\sqrt{899}}x\right).
\label{mab:rm:eq:chaos-profile}
\end{equation}

Le développement à l’origine est

\begin{equation}
u_0(x)=x^2-\frac{715769}{2424603}x^4+O(x^6).
\label{mab:rm:eq:chaos-taylor}
\end{equation}

La somme finie admet l’expression fermée suivante. Si \(y=x/\sqrt{899}\),

\begin{equation}
u_0(x)=1-\frac{\cos(37y)\sin(36y)}{12\sin(3y)}
=1-\frac{\sin(73y)-\sin y}{24\sin(3y)}.
\label{mab:rm:eq:chaos-closed}
\end{equation}

La singularité apparente en \(y=0\) est éliminable et reproduit \cref{mab:rm:eq:chaos-taylor}.

\begin{proof}[Dérivation de la forme fermée]
Les douze arguments forment une progression arithmétique :
\[
2(3m-1)y=(6m-2)y,\qquad m=1,\ldots,12.
\]
La somme des cosinus d’une progression satisfait
\[
\sum_{m=1}^{N}\cos(a+md)
=\frac{\sin(Nd/2)}{\sin(d/2)}
 \cos\!\left(a+\frac{N+1}{2}d\right).
\]
En prenant \(N=12\), \(d=6y\) et \(a=-2y\), on obtient
\[
\sum_{m=1}^{12}\cos((6m-2)y)
=\frac{\sin36y}{\sin3y}\cos37y.
\]
L’identité de transformation produit--en--somme
\(2\cos37y\sin36y=\sin73y-\sin y\) donne la deuxième expression de
\cref{mab:rm:eq:chaos-closed}.
\end{proof}

Le coefficient quartique n’est pas non plus ajusté. Les sommes entières
\begin{equation}
\sum_{m=1}^{12}(3m-1)^2=5394,\qquad
\sum_{m=1}^{12}(3m-1)^4=4294614
\label{mab:rm:eq:chaos-moments24}
\end{equation}
insérées dans le développement de \(\cos z\) produisent exactement
\[
\frac1{24\cdot12}\sum_m
\left(\frac{2(3m-1)}{\sqrt{899}}\right)^4
=\frac{715769}{2424603}.
\]
Ainsi, tant la normalisation quadratique que la première correction non
linéaire sont des moments finis de la même roue.

\section{Appartenance à la classe analytique unimodale quadratique}

Considérons

\[
f_\lambda(x)=1-\lambda u_0(x),\qquad \lambda>0.
\]

Pour \(0<x\le1\), tous les arguments

\[
\frac{2(3m-1)}{\sqrt{899}}x
\]

appartiennent à \((0,\pi)\). Par conséquent,

\[
u_0'(x)>0,\qquad u_0'''(x)<0.
\]

Comme \(u_0\) est paire, \(x=0\) est son unique point critique et \(u_0''(0)=2\). La dérivée schwarzienne de \(f_\lambda\) est

\begin{equation}
S(f_\lambda)
=\frac{u_0'''}{u_0'}
-\frac32\left(\frac{u_0''}{u_0'}\right)^2<0
\qquad(0<|x|\le1).
\label{mab:rm:eq:negative-schwarzian}
\end{equation}

\begin{theorem}[Appartenance exacte à la classe unimodale]
Pour
\[
0<\lambda\le \frac{2}{u_0(1)},
\]
la famille \(f_\lambda=1-\lambda u_0\) est une application réelle--analytique de \([-1,1]\) dans lui-même, paire, avec un unique point critique quadratique à l’origine et une dérivée schwarzienne négative en dehors de ce point. Elle appartient donc, sans utiliser les valeurs de Feigenbaum comme valeur cible, à la classe \(S\)-unimodale quadratique sur laquelle agit le renormalisateur de doublement de période \cite{rm:Feigenbaum1978,rm:Lanford1982}. Pour \(\lambda\) en dehors de cet intervalle, le germe analytique et sa forme locale sont conservés, mais l’application de l’intervalle dans lui-même n’est pas automatiquement affirmée.
\end{theorem}

Le théorème établit la contribution intrinsèque de HMT préalable au théorème
d’universalité : le germe analytique et son type de criticité.

\begin{proof}[Détails du signe et de l’application dans soi-même]
En dérivant \cref{mab:rm:eq:chaos-profile},
\begin{align*}
u_0'(x)&=\frac1{12}\sum_m k_m\sin(k_mx),\\
u_0'''(x)&=-\frac1{12}\sum_m k_m^3\sin(k_mx),
\qquad k_m=\frac{2(3m-1)}{\sqrt{899}}.
\end{align*}
Pour \(0<x\le1\), \(0<k_mx<\pi\), donc tous les termes de la première
ligne sont positifs et ceux de la seconde négatifs. La parité étend la
monotonie aux deux côtés de l’origine. Comme \(u_0(0)=0\) et
\(0\le u_0(x)\le u_0(1)\), la borne
\(0<\lambda\le2/u_0(1)\) implique
\(-1\le1-\lambda u_0(x)\le1\). Enfin,
\(u_0''(0)=1/12\sum k_m^2=2\), de sorte que le point critique est
quadratique.
\end{proof}

\section{Suites finies de paramètres superstables}

La suite superstable certifiée jusqu’au niveau huit produit

\[
\delta_8=4.66921218915\ldots,
\qquad
\alpha_8=2.50296847988\ldots.
\]

Son statut est celui de \emph{Certificat fini exempt de valeurs cibles} : chaque
approximant est obtenu à partir de la famille préalablement définie, sans introduire
comme donnée la limite universelle. La coïncidence décimale fournit un
indicateur de convergence, mais ne remplace pas le théorème d’hyperbolicité.

Une extension numérique indépendante de la même suite, avec un pas
adapté à la séparation entre les racines et un contrôle de la période exacte,
atteint provisoirement :

\begin{center}
\small
\begin{tabular}{@{}ccll@{}}
\toprule
Niveau & \(\lambda_n\) & \(\delta_n\) & \(\alpha_n\)\\
\midrule
11 & 1.87403828195439908045 & 4.66920170694424982745 & 2.50290847517165181764\\
12 & 1.87403836411023460451 & 4.66920163036308158848 & 2.50290800379001546915\\
13 & 1.87403838170549870372 & 4.66920161361839219913 & 2.50290790268748193108\\
14 & 1.87403838547386545431 & 4.66920161007472591023 & 2.50290788100969754274\\
\bottomrule
\end{tabular}
\end{center}

Ces niveaux ont le statut de \emph{Nouveau certificat numérique
en attente d’un enregistrement reproductible} ; ils ne remplacent pas le certificat persistant
jusqu’au niveau huit. Une extrapolation d’Aitken fournit
\(\lambda_\infty^{\HMT}\simeq1.8740383865008918080\) et
\(\alpha_F\simeq2.50290787509\) comme candidats à la limite, sous réserve de
la démonstration de l’hyperbolicité et de la transversalité.

Pour éviter une procédure numérique non explicitée, la méthode est formulée comme
suit. Définissez
\begin{equation}
F_n(\lambda)=f_\lambda^{\,2^n}(0).
\label{mab:rm:eq:superstable-equation}
\end{equation}
La racine superstable \(\lambda_n\) est la première racine de \(F_n\) à
droite de \(\lambda_{n-1}\) qui satisfait simultanément
\[
f_{\lambda_n}^{\,2^j}(0)\ne0
\quad(0\le j<n).
\]
Cette inégalité exclut les périodes diviseurs. Avec des intervalles disjoints
\(I_n\ni\lambda_n\), on calcule
\begin{equation}
\delta_n=
\frac{\lambda_{n-1}-\lambda_{n-2}}
     {\lambda_n-\lambda_{n-1}},
\qquad
\alpha_n=-\frac{d_{n-1}}{d_n},
\label{mab:rm:eq:feigenbaum-approximants}
\end{equation}
où \(d_n=f_{\lambda_n}^{\,2^{n-1}}(0)\) est la distance orientée du
dernier point nouveau au point critique. Un certificat persistant doit conserver
les intervalles de racine, l’exclusion des périodes inférieures et la propagation
de l’erreur de \cref{mab:rm:eq:feigenbaum-approximants} ; un simple nombre décimal imprimé ne
remplit pas ces trois conditions.

La règle de discrétisation constitue également un contrôle négatif. Une
borne inférieure fixe supérieure à la séparation
\(|\lambda_n-\lambda_{n-1}|\) peut manquer la première racine et produire
néanmoins une suite numériquement plausible. C’est pourquoi l’incrément
doit être proportionné à la séparation précédente et chaque candidat doit être vérifié
par sa période exacte.

\section{Conditions de certification de l’opérateur de renormalisation}

Sur un espace de fonctions paires analytiques, nous définissons

\begin{equation}
\mathcal R(f)(x)=-a(f)^{-1}f\bigl(f(-a(f)x)\bigr),
\qquad a(f)=-f(1).
\label{mab:rm:eq:renormalization}
\end{equation}

La clôture complète de \(\delta_F\) et de \(\alpha_F\) exige de certifier conjointement :

\begin{enumerate}
\item l’existence d’un point fixe \(g=\mathcal R(g)\) dans une boule analytique déclarée ;
\item l’unicité locale de ce point fixe ;
\item une unique valeur propre réelle instable de \(D\mathcal R_g\), égale à \(\delta_F\) ;
\item un rayon spectral strictement inférieur à un sur le complément stable ;
\item la transversalité de la courbe \(\lambda\mapsto f_\lambda\) par rapport à la variété stable de \(g\).
\end{enumerate}

Une implémentation HMT naturelle emploie des troncatures de Chebyshev d’ordre
\(9^m\), de l’arithmétique d’intervalles et des polynômes de rayons. La preuve doit
fournir une boule, un inverse approché, une borne du défaut et une borne
de Lipschitz satisfaisant \(Y+Z(r)<r\). Ainsi, la profondeur
nonadique paramètre un certificat de point fixe et non une simple suite
d’approximations décimales.

Dans une base paire de Chebyshev,
\[
g(x)=\sum_{j=0}^{N}g_{2j}T_{2j}(x),\qquad N=9^m,
\]
le résidu fini est \(F_N(g)=\Pi_N(\mathcal Rg-g)\). Si \(A_N\) approche
\((DF_N(\bar g))^{-1}\), les quantités à certifier sont
\begin{align*}
Y&=\|A_NF_N(\bar g)\|,\\
Z_0&=\|I-A_NDF_N(\bar g)\|,\\
Z_1(r)&=\sup_{\|h\|\le r}
\|A_N(DF_N(\bar g+h)-DF_N(\bar g))\|.
\end{align*}
Une inégalité
\(Y+(Z_0+Z_1(r))r<r\) enferme un unique point fixe. Le bloc spectral
ultérieur doit isoler une direction instable et contracter le complément. L’usage
de \(9^m\) intervient ainsi comme loi de raffinement du certificat, non comme
décoration numérologique.

\begin{corollary}[Confluence exacte du mode trente]
Le mode \(30\) admet trois réalisations simultanées et typologiquement
distinctes :
\begin{align*}
30&\longmapsto(3\cdot30,4\cdot30)=(90,120)
&&\text{dans la complétion constitutive},\\
30&\longmapsto30^2-1=899=29\cdot31
&&\text{dans la normalisation critique},\\
30&\longmapsto
M_{30}=\begin{pmatrix}30&899\\1&30\end{pmatrix}
&&\text{dans la transformation hyperbolique de Pell}.
\end{align*}
Les trois identités sont exactes et sont obtenues sans prescrire de valeurs cibles. L’invariant conservé est
l’entier central \(30\) avec son orientation : extension \(3\to4\),
défaut quadratique \(-1\) et unité de norme un. Ce qui n’est pas affirmé est
l’égalité entre l’opérateur de vide et le renormalisateur de doublement.
\end{corollary}

\begin{remark}[Contrôle négatif]
Une deuxième valeur propre en dehors du disque unité, une boule de rayons sans solution, une perte d’unicité ou la tangence de la famille HMT à la variété stable réfutent la promotion universelle. Aucune ne réfute le germe exact ni le théorème de la dérivée schwarzienne.
\end{remark}

\begin{proposition}[Résultat principal]
La signification structurelle de la constante de doublement réside dans
l’annulation dodécaphasique de \(\pi\), dans l’identité
\(899=30^2-1\), dans la forme trigonométrique fermée et dans l’appartenance
démontrée à la classe analytique sur laquelle agit le renormalisateur. Son
développement décimal constitue exclusivement l’évaluation terminale.
\end{proposition}

\subsection*{Portée logique et domaine de validité}
Profil, Taylor, normalisateur, unité de Pell, unimodalité et dérivée schwarzienne : \emph{Théorème interne exact}. Approximants jusqu’au niveau huit : \emph{Certificat fini persistant}. Niveaux 11--14 : \emph{Nouveau certificat numérique en attente de pérennisation}. Hyperbolicité et transversalité du renormalisateur : \emph{Programme ouvert avec critère de réfutation exécutable}.

